\hypertarget{appendix-1823a14330cb}{}
\section*{Interval arithmetic}
\chapter{Interval arithmetic}\label{appendix-1823a14330cb-interval-arithmetic}

\textbf{Curriculum depth:} Introduced · \textbf{First seen:} C004 · \textbf{Latest development:} C004

\section{Definition and mechanism}\label{appendix-1823a14330cb-definition-and-mechanism}

An interval encloses a possible exact value between lower and upper bounds. Correct outward rounding can preserve containment when the whole algorithm honors necessary assumptions.

\section{Prerequisites}\label{appendix-1823a14330cb-prerequisites}

\begin{itemize}
\tightlist
\item
  \hyperref[appendix-637d164ae5ca]{Directed rounding modes} --- The signed rounding direction is state that a demonstration or enclosure algorithm must control.
\item
  \hyperref[appendix-f11d1a521253]{Floating-point environment and compiler modes} --- Compiler/runtime assumptions must preserve the intended direction and special-value semantics.
\end{itemize}

\section{Used by}\label{appendix-1823a14330cb-used-by}

Connect later dependent concepts here as their lessons arrive.

\section{Recall and next depth}\label{appendix-1823a14330cb-recall-and-next-depth}

Reconstruct the concrete example in the definition and identify its validity assumptions.

\textbf{Pending:} Directed rounding alone does not establish a valid numerical proof or robust predicate.

\section{Source and navigation}\label{appendix-1823a14330cb-source-and-navigation}

\href{Continuous\%20Notes/Computing\%20Foundations\%20-\%20From\%20Bits\%20to\%20C++\%20and\%20Systems\%20APIs\#A\%20stretching\%20ruler\%20for\%20real-valued\%20quantities}{M004 teaching chapter} · \hyperref[appendix-fd090d982a15-m004---parts-61-76]{Per-part definitions} · \href{../02-Foundations-II/main.pdf\#nameddest=appendix-4ad47bdd269a}{Dependency Map (02-Foundations-II, appendix-4ad47bdd269a)}
